\chapter{2007年大纲II卷（文）}

\section{单选题}

\begin{problem}
\(\cos 330^{\circ}=\)（\quad）

\choices
    {\(\frac{1}{2}\)}
    {\(-\frac{1}{2}\)}
    {\(\frac{\sqrt{3}}{2}\)}
    {\(-\frac{\sqrt{3}}{2}\)}
\end{problem}

\begin{problem}
设集合 \(U=\{1,2,3,4\}\)，\(A=\{1,2\}\)，\(B=\{2,4\}\)，则 \(\complement_U(A\cup B)=\)（\quad）

\choices
    {\(\{2\}\)}
    {\(\{3\}\)}
    {\(\{1,2,4\}\)}
    {\(\{1,4\}\)}
\end{problem}

\begin{problem}
函数 \(f(x)=|\sin x|\) 的一个单调递增区间是（\quad）

\choices
    {\(\left(-\frac{\pi}{4},\frac{\pi}{4}\right)\)}
    {\(\left(\frac{\pi}{4},\frac{3\pi}{4}\right)\)}
    {\(\left(\pi,\frac{3\pi}{2}\right)\)}
    {\(\left(\frac{3\pi}{2},2\pi\right)\)}
\end{problem}

\begin{problem}
以下四个数中的最大者是（\quad）

\choices
    {\((\ln 2)^2\)}
    {\(\ln\left(\ln 2\right)\)}
    {\(\ln\sqrt{2}\)}
    {\(\ln 2\)}
\end{problem}

\begin{problem}
不等式 \(\frac{x-2}{x+3}>0\) 的解集是（\quad）

\choices
    {\(\left(-3,2\right)\)}
    {\(\left(2,+\infty\right)\)}
    {\(\left(-\infty,-3\right)\cup\left(2,+\infty\right)\)}
    {\(\left(-\infty,-2\right)\cup\left(3,+\infty\right)\)}
\end{problem}

\begin{problem}
在 \(\triangle ABC\) 中，已知 \(D\) 是 \(AB\) 边上一点，若 \(\overrightarrow{AD}=2\overrightarrow{DB}\)，\(\overrightarrow{CD}=\frac{1}{3}\overrightarrow{CA}+\lambda\overrightarrow{CB}\)，则 \(\lambda=\)（\quad）

\choices
    {\(\frac{2}{3}\)}
    {\(\frac{1}{3}\)}
    {\(-\frac{1}{3}\)}
    {\(-\frac{2}{3}\)}
\end{problem}

\begin{problem}
已知正三棱锥的侧棱长与底面边长的 \(2\) 倍，则侧棱与底面所成角的余弦值等于（\quad）

\choices
    {\(\frac{\sqrt{3}}{6}\)}
    {\(\frac{\sqrt{3}}{4}\)}
    {\(\frac{\sqrt{2}}{2}\)}
    {\(\frac{\sqrt{3}}{2}\)}
\end{problem}

\begin{problem}
已知曲线 \(y=\frac{x^2}{4}\) 的一条切线的斜率为 \(\frac{1}{2}\)，则切点的横坐标为（\quad）

\choices
    {\(1\)}
    {\(2\)}
    {\(3\)}
    {\(4\)}
\end{problem}

\begin{problem}
把函数 \(y=\e^x\) 的图象按向量 \(\bs{a}=(2,0)\) 平移，得到 \(y=f(x)\) 的图象，则 \(f(x)=\)（\quad）

\choices
    {\(\e^x+2\)}
    {\(\e^x-2\)}
    {\(\e^{x-2}\)}
    {\(\e^{x+2}\)}
\end{problem}

\begin{problem}
\(5\) 位同学报名参加两个课外活动小组，每位同学限报其中的一个小组，则不同的报名方法共有（\quad）

\choices
    {\(10\text{ 种}\)}
    {\(20\text{ 种}\)}
    {\(25\text{ 种}\)}
    {\(32\text{ 种}\)}
\end{problem}

\begin{problem}
已知椭圆的长轴长是短轴长的 \(2\) 倍，则椭圆的离心率为（\quad）

\choices
    {\(\frac{1}{3}\)}
    {\(\frac{\sqrt{3}}{3}\)}
    {\(\frac{1}{2}\)}
    {\(\frac{\sqrt{3}}{2}\)}
\end{problem}

\begin{problem}
设 \(F_1\)，\(F_2\) 分别是双曲线 \(x^2-\frac{y^2}{9}=1\) 的左，右焦点，若点 \(P\) 在双曲线上，且 \(\overrightarrow{PF_1}\cdot\overrightarrow{PF_2}=0\)，则 \(\left|\overrightarrow{PF_1}+\overrightarrow{PF_2}\right|=\)（\quad）

\choices
    {\(\sqrt{10}\)}
    {\(2\sqrt{10}\)}
    {\(\sqrt{5}\)}
    {\(2\sqrt{5}\)}
\end{problem}

\section{填空题}

\begin{problem}
一个总体含有 \(100\) 个个体，以简单随机抽样方式从该总体中抽取一个容量为 \(5\) 的样本，则指定的某个个体被抽到的概率为\fillinblank{}.
\end{problem}

\begin{problem}
已知数列的通项 \(a_n=-5n+2\)，则其前 \(n\) 项和为 \(S_n=\)\fillinblank{}.
\end{problem}

\begin{problem}
一个正四棱柱的各个顶点在一个直径为 \(2\text{ cm}\) 的球面上. 如果正四棱柱的底面边长为 \(1\text{ cm}\)，那么该棱柱的表面积为\fillinblank{}\(\text{cm}^2\).
\end{problem}

\begin{problem}
\((1+2x^2)\left(x+\frac{1}{x}\right)^8\) 的展开式中常数项为\fillinblank{}.（用数字作答）
\end{problem}

\section{解答题}

\begin{problem}
设等比数列 \(\{a_n\}\) 的公比 \(q<1\)，前 \(n\) 项和为 \(S_n\)，已知 \(a_3=2\)，\(S_4=5S_2\)，求 \(\{a_n\}\) 的通项公式.
\end{problem}

\begin{problem}
在 \(\triangle ABC\) 中，已知内角 \(A=\frac{\pi}{3}\)，边 \(BC=2\sqrt{3}\)，设内角 \(B=x\)，周长为 \(y\).

\begin{enumerate}
\item 求函数 \(y=f(x)\) 的解析式和定义域；
\item 求 \(y\) 的最大值.
\end{enumerate}
\end{problem}

\begin{problem}
从某批产品中，有放回地抽取产品二次，每次随机抽取 \(1\) 件，假设事件 \(A\)：“取出的 \(2\) 件产品中至多有 \(1\) 件是二等品”的概率 \(P(A)=0.96\).

\begin{enumerate}
\item 求从该批产品中任取 \(1\) 件是二等品的概率 \(p\)；
\item 若该批产品共有 \(100\) 件，从中任意抽取 \(2\) 件，求事件 \(B\)：“取出的 \(2\) 件产品中至少有一件二等品”的概率 \(P(B)\).
\end{enumerate}
\end{problem}

\begin{problem}
如图，在四棱锥 \(S-ABCD\) 中，底面 \(ABCD\) 为正方形，侧棱 \(SD\perp\) 底面 \(ABCD\)，\(E\)，\(F\) 分别是 \(AB\)，\(SC\) 的中点.

\begin{enumerate}
\item 求证：\(EF\parallel\) 平面 \(SAD\)；
\item 设 \(SD=2CD\)，求二面角 \(A-EF-D\) 的大小.
\end{enumerate}

\bitmapfigure[float=inline,width=0.38\linewidth]{img/2007/syllabus_paper_2_liberal/q20_fig1.png}
\end{problem}

\begin{problem}
在直角坐标系 \(xOy\) 中，以 \(O\) 为圆心的圆与直线 \(x-\sqrt{3}y=4\) 相切.

\begin{enumerate}
\item 求圆 \(O\) 的方程；
\item 圆 \(O\) 与 \(x\) 轴相交于 \(A\)，\(B\) 两点，圆内的动点 \(P\) 使 \(|PA|\)，\(|PO|\)，\(|PB|\) 成等比数列，求 \(\overrightarrow{PA}\cdot\overrightarrow{PB}\) 的取值范围.
\end{enumerate}
\end{problem}

\begin{problem}
已知函数 \(f(x)=\frac{1}{3}ax^3-bx^2+(2-b)x+1\) 在 \(x=x_1\) 处取得极大值，在 \(x=x_2\) 处取得极小值，且 \(0<x_1<1<x_2<2\).

\begin{enumerate}
\item 证明 \(a>0\)；
\item 若 \(z=a+2b\)，求 \(z\) 的取值范围.
\end{enumerate}
\end{problem}
